\documentclass[12pt]{article}
\usepackage{fullpage}
\usepackage[T1]{fontenc}
\usepackage[utf8]{inputenc}
\usepackage{lmodern}
\usepackage{microtype}
\usepackage{amsmath,amssymb}
\usepackage{graphicx}
\usepackage{booktabs}
\usepackage{hyperref}
\usepackage{url}
\usepackage{color}
\usepackage{xcolor}
\usepackage{enumitem}
\usepackage{amsfonts}
\usepackage{amsthm}
\usepackage{algorithm}
\usepackage{algpseudocode}

\newtheorem{theorem}{Theorem}
\newtheorem{proposition}[theorem]{Proposition}
\newtheorem{lemma}[theorem]{Lemma}
\newtheorem{corollary}[theorem]{Corollary}
\theoremstyle{definition}
\newtheorem{definition}[theorem]{Definition}
\newtheorem{remark}[theorem]{Remark}

\DeclareMathOperator{\vecop}{vec}
\DeclareMathOperator{\diag}{diag}
\DeclareMathOperator{\mat}{mat}

\title{Preconditioned Conjugate Gradient for RKHS-Constrained CP Tensor Decomposition}
\author{}
\date{}

\begin{document}
\maketitle

\section{Problem Setup and Notation}

We work with a $d$-way tensor $\mathcal{T} \in \mathbb{R}^{n_1 \times \cdots \times n_d}$ with
missing entries.  For the mode-$k$ alternating-least-squares subproblem we fix all factor
matrices $A_1,\dots,A_{k-1},A_{k+1},\dots,A_d$ and solve for the RKHS-constrained factor
$A_k = KW$, where $K \in \mathbb{R}^{n \times n}$ ($n \equiv n_k$) is a symmetric positive
semidefinite (PSD) kernel matrix and $W \in \mathbb{R}^{n \times r}$ is the unknown.

The key quantities are:
\begin{itemize}
  \item $N = \prod_i n_i$ (total number of tensor entries), $n \equiv n_k$,
        $M = \prod_{i \neq k} n_i$ (so $N = nM$).
  \item $q \ll N$: the number of observed entries; $n, r < q$.
  \item $T \in \mathbb{R}^{n \times M}$: mode-$k$ unfolding with zeros at unobserved positions.
  \item $S \in \mathbb{R}^{N \times q}$: selection matrix (columns are standard basis vectors
        $e_\omega$, one per observed entry $\omega$); $S^T \vecop(T)$ extracts the $q$ observed
        values.
  \item $Z = A_d \odot \cdots \odot A_{k+1} \odot A_{k-1} \odot \cdots \odot A_1
        \in \mathbb{R}^{M \times r}$: Khatri-Rao product of all non-$k$ factor matrices.
  \item $B = TZ \in \mathbb{R}^{n \times r}$: the Matricized Tensor Times Khatri-Rao Product
        (MTTKRP).
  \item $\Omega \subseteq [n] \times [M]$: the set of observed index pairs
        $\omega = (i_\omega, j_\omega)$ where $i_\omega \in [n]$ and $j_\omega \in [M]$,
        with $|\Omega| = q$.
\end{itemize}

The subproblem to be solved is the linear system
\begin{equation}
  \mathbf{H}\,\vecop(W) = (I_r \otimes K)\,\vecop(B),
  \label{eq:system}
\end{equation}
where
\begin{equation}
  \mathbf{H} = (Z \otimes K)^T S S^T (Z \otimes K) + \lambda (I_r \otimes K),
  \label{eq:H}
\end{equation}
$\lambda > 0$ is a regularization parameter, and $I_r$ is the $r \times r$ identity.
The system size is $nr \times nr$.

\paragraph{Compact notation for observed rows.}
We write $\widetilde{K} \in \mathbb{R}^{q \times n}$ for the matrix whose $\omega$-th row is
$K_{i_\omega,:}$ (the $i_\omega$-th row of $K$), and $\widetilde{Z} \in \mathbb{R}^{q \times r}$
for the matrix whose $\omega$-th row is $Z_{j_\omega,:}$ (the $j_\omega$-th row of $Z$).
Both $\widetilde{K}$ and $\widetilde{Z}$ are formed by selecting $q$ rows from $K$ and $Z$
respectively, with possible repetition.

\section{Structure of the System Matrix $\mathbf{H}$}

\subsection{Symmetry and Positive Definiteness}

\begin{proposition}
  The matrix $\mathbf{H}$ defined in \eqref{eq:H} is symmetric positive definite whenever $K$
  is positive definite and $\lambda > 0$.
\end{proposition}

\begin{proof}
  The term $(Z \otimes K)^T SS^T (Z \otimes K)$ is symmetric positive semidefinite since
  $SS^T$ is symmetric PSD (it is an orthogonal projection onto the $q$-dimensional subspace
  spanned by observed entries) and the Kronecker product $(Z \otimes K)$ has full column rank
  when $K$ is invertible.  The term $\lambda(I_r \otimes K)$ is symmetric positive definite
  when $\lambda > 0$ and $K \succ 0$.  Their sum is therefore symmetric positive definite.
\end{proof}

\subsection{Entry-Wise Decomposition}

We expand the data-fidelity term using the structure of $S$.  Since $SS^T$ is the $N \times N$
diagonal matrix with ones at positions corresponding to observed entries, we have
\[
  (Z \otimes K)^T SS^T (Z \otimes K)
  = \sum_{\omega \in \Omega}
    \bigl(Z_{j_\omega,:} \otimes K_{i_\omega,:}\bigr)^T
    \bigl(Z_{j_\omega,:} \otimes K_{i_\omega,:}\bigr).
\]
Using the identity $(a \otimes b)^T(a \otimes b) = (aa^T) \otimes (bb^T)$ for row vectors:
\begin{equation}
  \mathbf{H}
  = \sum_{\omega \in \Omega}
    \bigl(Z_{j_\omega,:}^T Z_{j_\omega,:}\bigr) \otimes
    \bigl(K_{i_\omega,:}^T K_{i_\omega,:}\bigr)
    + \lambda I_r \otimes K.
  \label{eq:H-sum}
\end{equation}
This sum has $q$ rank-1 terms (in the Kronecker sense), each of size $nr \times nr$.
Explicitly forming $\mathbf{H}$ from \eqref{eq:H-sum} costs $O(q n^2 r^2)$, which is
prohibitive.  Solving the explicitly formed system by Cholesky factorization costs $O(n^3 r^3)$.
Both computations are avoided by PCG.

\section{Direct Symmetric Method}

For small problems (modest $n$ and $r$) one can assemble $\mathbf{H}$ and solve directly.

\begin{algorithm}
\caption{Direct Symmetric Solver for the RKHS-CP Mode-$k$ Subproblem}
\label{alg:direct}
\begin{algorithmic}[1]
  \Require Kernel $K \in \mathbb{R}^{n \times n}$,
           Khatri-Rao product $Z \in \mathbb{R}^{M \times r}$,
           MTTKRP $B \in \mathbb{R}^{n \times r}$,
           observed index set $\Omega$ (pairs $(i_\omega,j_\omega)$),
           regularization $\lambda > 0$.
  \Ensure  $W \in \mathbb{R}^{n \times r}$ solving \eqref{eq:system}.

  \State \textbf{[Assemble $\mathbf{H}$]} Initialize $\mathbf{H} \leftarrow \mathbf{0} \in \mathbb{R}^{nr \times nr}$.
  \For{each $\omega = (i_\omega, j_\omega) \in \Omega$}
    \State $\mathbf{H} \leftarrow \mathbf{H}
              + (Z_{j_\omega,:}^T Z_{j_\omega,:}) \otimes (K_{i_\omega,:}^T K_{i_\omega,:})$
    \Comment{rank-1 Kronecker update; cost $O(n^2 r^2)$ per step}
  \EndFor
  \State $\mathbf{H} \leftarrow \mathbf{H} + \lambda (I_r \otimes K)$
         \Comment{add regularizer; cost $O(n^2 r)$}

  \State \textbf{[Form RHS]} $b \leftarrow \vecop\!\bigl((I_r \otimes K)\vecop(B)\bigr)
                                         = \vecop(KB)$
         \Comment{cost $O(n^2 r)$}

  \State \textbf{[Cholesky]} $\mathbf{H} = LL^T$
         \Comment{cost $O(n^3 r^3)$}

  \State Solve $L y = b$, then $L^T \vecop(W) = y$
         \Comment{cost $O(n^2 r^2)$}

  \State \Return $W = \mat_{n \times r}(\vecop(W))$
\end{algorithmic}
\end{algorithm}

\paragraph{Complexity.}
Assembling $\mathbf{H}$ costs $O(q n^2 r^2)$; the Cholesky factorization costs $O(n^3 r^3)$;
the two triangular solves cost $O(n^2 r^2)$.  The dominant term is $O(q n^2 r^2 + n^3 r^3)$.

\section{Efficient Matrix-Vector Product for PCG}

PCG requires only \emph{matrix-vector products} $v \mapsto \mathbf{H} v$ (no explicit matrix
formation).  We show how to compute $\mathbf{H} v$ in $O(n^2 r + qr)$ time, avoiding any
$O(N)$ or $O(M)$ operations.

Let $v = \vecop(V)$ with $V \in \mathbb{R}^{n \times r}$.  Split the product:
\begin{equation}
  \mathbf{H}\,\vecop(V)
  = \underbrace{(Z \otimes K)^T SS^T (Z \otimes K)\,\vecop(V)}_{\text{data term}}
    + \underbrace{\lambda(I_r \otimes K)\,\vecop(V)}_{\text{regularizer term}}.
  \label{eq:Hv-split}
\end{equation}

\subsection{Regularizer Term}

Using the Kronecker-vec identity $(A \otimes B)\vecop(X) = \vecop(BXA^T)$:
\[
  \lambda(I_r \otimes K)\,\vecop(V) = \lambda\,\vecop(KV I_r^T) = \lambda\,\vecop(KV).
\]
\textbf{Cost:} Compute $KV \in \mathbb{R}^{n \times r}$ as $r$ matrix-vector products of the
$n \times n$ matrix $K$: $O(n^2 r)$.

\subsection{Data Term}

We factor the data term into three sequential operations.

\paragraph{Step~A: Apply $(Z \otimes K)$ to $\vecop(V)$.}

We need only the $q$ entries of the full $N$-vector $(Z \otimes K)\vecop(V)$ that are selected
by $S^T$.  For observed entry $\omega = (i_\omega, j_\omega)$:
\begin{align}
  \bigl[(Z \otimes K)\,\vecop(V)\bigr]_\omega
  &= \bigl(Z_{j_\omega,:} \otimes K_{i_\omega,:}\bigr)\,\vecop(V) \notag \\
  &= K_{i_\omega,:}\,V\,Z_{j_\omega,:}^T
  \label{eq:u-omega}
\end{align}
where we have used the bilinear form identity
$(a^T \otimes b^T)\vecop(V) = a^T V^T b$ (applied after recognizing that
$(Z_{j_\omega,:} \otimes K_{i_\omega,:})\vecop(V) = K_{i_\omega,:} V Z_{j_\omega,:}^T$
is a scalar, the trace $\mathrm{tr}(K_{i_\omega,:}^T Z_{j_\omega,:} \cdot V^T)$).
More explicitly:
\[
  u_\omega := K_{i_\omega,:}\,V\,Z_{j_\omega,:}^T
            = \sum_{m=1}^n \sum_{s=1}^r K_{i_\omega, m}\,V_{ms}\,Z_{j_\omega, s} \in \mathbb{R}.
\]

\textbf{Efficient computation.}  Pre-compute $\widetilde{V} = KV \in \mathbb{R}^{n \times r}$
at cost $O(n^2 r)$.  Then for each $\omega \in \Omega$:
\[
  u_\omega = \widetilde{V}_{i_\omega,:} \cdot Z_{j_\omega,:}^T
           = \widetilde{K}_{i_\omega,:}\,V\,\widetilde{Z}_{j_\omega,:}^T,
\]
i.e., the inner product of the $i_\omega$-th row of $\widetilde{V}$ with the $j_\omega$-th row
of $Z$ (equivalently, the $\omega$-th row of $\widetilde{Z}$).  Computing all $q$ scalars
$\{u_\omega\}_{\omega \in \Omega}$ requires one inner product of length $r$ per observation:
\textbf{cost} $O(qr)$.

In matrix form, $u = \diag(\widetilde{K} V \widetilde{Z}^T) \in \mathbb{R}^q$, where the
diagonal picks the matching row indices.  More precisely, with $\widetilde{V} := KV$,
\[
  u_\omega = \bigl(\widetilde{V}\bigr)_{i_\omega,:}\,\widetilde{Z}_{\omega,:}^T,
  \quad \omega = 1, \dots, q,
\]
which is the elementwise product summed over $r$ columns:
$u = \mathrm{rowsum}\bigl(\widetilde{V}[\text{row-idx}] \circ \widetilde{Z}\bigr)$,
where $\widetilde{V}[\text{row-idx}]$ selects row $i_\omega$ of $\widetilde{V}$ for each
$\omega$ (a $q \times r$ gather), $\circ$ denotes elementwise product, and $\mathrm{rowsum}$
sums across the $r$ columns.  This is a standard sparse-dense gather-and-dot, cost $O(qr)$.

\paragraph{Step~B: Apply $SS^T$.}

Since $SS^T$ is the projection onto observed entries, it maps a full $N$-vector to itself with
all unobserved coordinates zeroed.  Because we computed only $\{u_\omega\}$ in Step~A, this
step is trivial: the output vector restricted to observed coordinates is $\{u_\omega\}$, and
all other entries are zero.  \textbf{Cost:} $O(1)$ (no work beyond the bookkeeping in Step~A).

\paragraph{Step~C: Apply $(Z \otimes K)^T$ to $SS^T u$.}

The full $N$-vector $SS^T u$ is supported on $\Omega$ with values $\{u_\omega\}$.  Therefore:
\begin{equation}
  (Z \otimes K)^T (SS^T u)
  = \sum_{\omega \in \Omega} u_\omega
    \bigl(Z_{j_\omega,:} \otimes K_{i_\omega,:}\bigr)^T,
  \label{eq:step-c}
\end{equation}
where each $(Z_{j_\omega,:} \otimes K_{i_\omega,:})^T \in \mathbb{R}^{nr}$ is reshaped as an
$n \times r$ matrix $K_{:,i_\omega} Z_{j_\omega,:}$ (outer product of the $i_\omega$-th column
of $K$ with the $j_\omega$-th row of $Z$).  Thus, reshaping the result to $\mathbb{R}^{n \times r}$:
\[
  \mat_{n \times r}\!\bigl[(Z \otimes K)^T (SS^T u)\bigr]
  = \sum_{\omega \in \Omega} u_\omega\,K_{:,i_\omega}\,Z_{j_\omega,:}
  = K \cdot G,
\]
where $G \in \mathbb{R}^{n \times r}$ is defined by
\begin{equation}
  G_{m,:} = \sum_{\omega \in \Omega:\, i_\omega = m} u_\omega\, Z_{j_\omega,:},
  \quad m = 1, \dots, n.
  \label{eq:G}
\end{equation}
That is, $G$ accumulates, for each mode-$k$ index $m$, the weighted sum of rows of $Z$
at observed entries with that mode-$k$ index.

\textbf{Efficient computation.}
$G$ is assembled from $q$ rank-1 (scalar times row-vector) contributions:
$G \leftarrow G + u_\omega e_{i_\omega} Z_{j_\omega,:}$ for each $\omega$, using only the
scalar $u_\omega$ and the row $Z_{j_\omega,:}$ (length $r$).  \textbf{Cost:} $O(qr)$.
Then $KG$ costs $O(n^2 r)$.

\subsection{Summary of Matrix-Vector Product}

\begin{algorithm}
\caption{Matrix-Vector Product $p \leftarrow \mathbf{H}\,\vecop(V)$}
\label{alg:matvec}
\begin{algorithmic}[1]
  \Require $V \in \mathbb{R}^{n \times r}$, kernel $K$, $Z$, observed set $\Omega$,
           parameter $\lambda$.
  \Ensure  $P \in \mathbb{R}^{n \times r}$ with $\vecop(P) = \mathbf{H}\,\vecop(V)$.

  \State $\widetilde{V} \leftarrow KV \in \mathbb{R}^{n \times r}$
         \Comment{cost $O(n^2 r)$}

  \State Initialize $G \leftarrow \mathbf{0} \in \mathbb{R}^{n \times r}$

  \For{each $\omega = (i_\omega, j_\omega) \in \Omega$}
    \State $u_\omega \leftarrow \widetilde{V}_{i_\omega,:} \cdot Z_{j_\omega,:}^T$
           \Comment{inner product of length $r$; cost $O(r)$}
    \State $G_{i_\omega,:} \leftarrow G_{i_\omega,:} + u_\omega \cdot Z_{j_\omega,:}$
           \Comment{accumulate; cost $O(r)$}
  \EndFor  \Comment{total loop cost: $O(qr)$}

  \State $P \leftarrow K G + \lambda K V = K(G + \lambda V)$
         \Comment{one $n \times n$ times $n \times r$ product; cost $O(n^2 r)$}

  \State \Return $P$
\end{algorithmic}
\end{algorithm}

\noindent
\textbf{Total cost per matrix-vector product:} $O(n^2 r + qr)$.

Note that the final step combines the data term $KG$ and the regularizer term $\lambda KV$ into
a single matrix product $K(G + \lambda V)$, saving one $O(n^2 r)$ multiplication.  Both $G$
and $V$ are $n \times r$ matrices, so their sum is formed at cost $O(nr)$ before the single
kernel multiplication.

\begin{remark}
  No operation in Algorithm~\ref{alg:matvec} involves vectors or matrices of size $N$ or $M$.
  Step~1 is $O(n^2 r)$ (kernel times an $n \times r$ matrix).  The loop visits each of the
  $q$ observed entries exactly once, performing $O(r)$ work per entry.  Step~5 is another
  $O(n^2 r)$ operation.  Since $n, r < q \ll N$ and $n \ll M$, all costs satisfy the
  constraint of avoiding $O(N)$ and $O(M)$ computations.
\end{remark}

\section{Preconditioner}

\subsection{Choice of Preconditioner}

We use the regularizer block as the preconditioner:
\begin{equation}
  P_{\mathrm{PCG}} = \lambda(I_r \otimes K).
  \label{eq:precond}
\end{equation}

\paragraph{Motivation.}
The matrix $\mathbf{H} = \mathcal{D} + \lambda(I_r \otimes K)$ where
$\mathcal{D} = (Z \otimes K)^T SS^T (Z \otimes K) \succeq 0$.  If the data term $\mathcal{D}$
is small relative to $\lambda(I_r \otimes K)$ (strong regularization or sparse observations),
then $\mathbf{H} \approx \lambda(I_r \otimes K)$ and the preconditioned system
$P_{\mathrm{PCG}}^{-1} \mathbf{H}$ has condition number close to~1.  More precisely, the
eigenvalues of $P_{\mathrm{PCG}}^{-1}\mathbf{H}$ lie in the interval
$[1,\, 1 + \sigma_{\max}(\mathcal{D})/\lambda\sigma_{\min}(K)]$, so the preconditioner
is effective when the ratio $\sigma_{\max}(\mathcal{D})/(\lambda\sigma_{\min}(K))$ is moderate.

\subsection{Applying the Preconditioner}

Solving $P_{\mathrm{PCG}}\,y = z$ for $y$ means solving
$\lambda(I_r \otimes K)\,\vecop(Y) = \vecop(Z_r)$, i.e.,
\[
  KY = \frac{1}{\lambda} Z_r,
\]
where $Z_r = \mat_{n \times r}(\vecop(z))$.  This requires solving $r$ linear systems with the
same matrix $K$.  We precompute the Cholesky factorization $K = LL^T$ once at cost $O(n^3)$
and then perform $r$ forward-backward substitutions at cost $O(n^2)$ each.

\textbf{Preconditioner application cost:} $O(n^3)$ (one-time Cholesky) $+ O(n^2 r)$ per PCG
iteration.

\section{Preconditioned Conjugate Gradient Algorithm}

\begin{algorithm}
\caption{PCG for RKHS-CP Mode-$k$ Subproblem}
\label{alg:pcg}
\begin{algorithmic}[1]
  \Require Kernel $K \in \mathbb{R}^{n \times n}$ (PSD),
           Khatri-Rao product $Z \in \mathbb{R}^{M \times r}$,
           MTTKRP $B \in \mathbb{R}^{n \times r}$,
           observed index set $\Omega$,
           regularization $\lambda > 0$,
           tolerance $\varepsilon > 0$,
           max iterations $t_{\max}$.
  \Ensure  Approximate solution $W \in \mathbb{R}^{n \times r}$ to \eqref{eq:system}.

  \State \textbf{[Preprocessing]}
         Compute Cholesky: $K = LL^T$
         \Comment{cost $O(n^3)$; done once}

  \State Compute $\widetilde{Z} \in \mathbb{R}^{q \times r}$:
         row $\omega$ is $Z_{j_\omega,:}$
         \Comment{cost $O(qr)$; sparse gather}

  \State Compute $\widetilde{K} \in \mathbb{R}^{q \times n}$:
         row $\omega$ is $K_{i_\omega,:}$
         \Comment{cost $O(qn)$; sparse gather}

  \State \textbf{[Initialize]} $W \leftarrow \mathbf{0} \in \mathbb{R}^{n \times r}$

  \State Form RHS: $b \leftarrow KB$
         \Comment{$\vecop(b) = (I_r \otimes K)\vecop(B)$; cost $O(n^2 r)$}

  \State $R \leftarrow b - \mathbf{H}(W)$
         \Comment{initial residual via Algorithm~\ref{alg:matvec}; cost $O(n^2 r + qr)$}

  \State Solve $K Y_0 = (1/\lambda) R$ for $Y_0$
         \Comment{preconditioner: $P_{\mathrm{PCG}}^{-1} R$; cost $O(n^2 r)$ using $L,L^T$}

  \State $D \leftarrow Y_0$
         \Comment{initial search direction}

  \State $\delta_{\mathrm{new}} \leftarrow \langle R, Y_0 \rangle_F$
         \Comment{$= \mathrm{tr}(R^T Y_0)$; cost $O(nr)$}

  \State $\delta_0 \leftarrow \delta_{\mathrm{new}}$

  \For{$t = 1, 2, \dots, t_{\max}$}

    \If{$\delta_{\mathrm{new}} \leq \varepsilon^2\,\delta_0$}
      \State \textbf{break}
    \EndIf

    \State $Q \leftarrow \mathbf{H}(D)$
           \Comment{mat-vec via Algorithm~\ref{alg:matvec}; cost $O(n^2 r + qr)$}

    \State $\alpha \leftarrow \delta_{\mathrm{new}} / \langle D, Q \rangle_F$
           \Comment{step size; cost $O(nr)$}

    \State $W \leftarrow W + \alpha D$
           \Comment{cost $O(nr)$}

    \State $R \leftarrow R - \alpha Q$
           \Comment{cost $O(nr)$}

    \State Solve $K Y = (1/\lambda) R$ for $Y$
           \Comment{preconditioner apply; cost $O(n^2 r)$}

    \State $\delta_{\mathrm{old}} \leftarrow \delta_{\mathrm{new}}$

    \State $\delta_{\mathrm{new}} \leftarrow \langle R, Y \rangle_F$
           \Comment{cost $O(nr)$}

    \State $\beta \leftarrow \delta_{\mathrm{new}} / \delta_{\mathrm{old}}$

    \State $D \leftarrow Y + \beta D$
           \Comment{update search direction; cost $O(nr)$}

  \EndFor

  \State \Return $W$
\end{algorithmic}
\end{algorithm}

\subsection{Correctness of PCG}

Algorithm~\ref{alg:pcg} is the standard preconditioned conjugate gradient method applied to the
symmetric positive definite system $\mathbf{H}\,\vecop(W) = \vecop(KB)$, with preconditioner
$P_{\mathrm{PCG}} = \lambda(I_r \otimes K)$.  The inner products $\langle \cdot, \cdot \rangle_F$
denote the Frobenius inner product on $\mathbb{R}^{n \times r}$, which is equivalent to the
Euclidean inner product on $\mathbb{R}^{nr}$ under the $\vecop$ identification.  Convergence in
exact arithmetic is guaranteed in at most $nr$ iterations; in practice, far fewer iterations
suffice because the preconditioner reduces the effective condition number.

\section{Complexity Analysis}

\subsection{Direct Method (Algorithm~\ref{alg:direct})}

\begin{center}
\begin{tabular}{lll}
  \toprule
  Step & Operation & Cost \\
  \midrule
  Assemble $\mathbf{H}$ & $q$ Kronecker rank-1 updates in $\mathbb{R}^{nr \times nr}$ & $O(q n^2 r^2)$ \\
  Add regularizer & $\lambda(I_r \otimes K)$ to $\mathbf{H}$ & $O(n^2 r)$ \\
  Form RHS & $KB$ & $O(n^2 r)$ \\
  Cholesky & $\mathbf{H} = LL^T$, $\mathbf{H} \in \mathbb{R}^{nr \times nr}$ & $O(n^3 r^3)$ \\
  Triangular solves & forward and backward substitution & $O(n^2 r^2)$ \\
  \midrule
  \textbf{Total} & & $O(q n^2 r^2 + n^3 r^3)$ \\
  \bottomrule
\end{tabular}
\end{center}

\subsection{PCG (Algorithm~\ref{alg:pcg})}

\begin{center}
\begin{tabular}{lll}
  \toprule
  Step & Operation & Cost \\
  \midrule
  Preprocessing & Cholesky of $K$; gather $\widetilde{Z}$, $\widetilde{K}$ & $O(n^3 + qn + qr)$ \\
  Per iteration (mat-vec) & Algorithm~\ref{alg:matvec} & $O(n^2 r + qr)$ \\
  Per iteration (precond) & $r$ triangular solves with $K$ & $O(n^2 r)$ \\
  Per iteration (other) & dot products, updates & $O(nr)$ \\
  \midrule
  Per iteration total & & $O(n^2 r + qr)$ \\
  \midrule
  \textbf{Total ($t$ iterations)} & & $O(n^3 + t(n^2 r + qr))$ \\
  \bottomrule
\end{tabular}
\end{center}

\paragraph{Iteration count and overall cost.}
Let $\kappa(P_{\mathrm{PCG}}^{-1}\mathbf{H})$ denote the condition number of the preconditioned
system.  PCG converges to $\varepsilon$-relative error in at most
$O(\sqrt{\kappa}\,\log(1/\varepsilon))$ iterations.  With the preconditioner
$P_{\mathrm{PCG}} = \lambda(I_r \otimes K)$, one can show
\[
  \kappa(P_{\mathrm{PCG}}^{-1}\mathbf{H})
  = \frac{\lambda_{\max}(\mathbf{H})}{\lambda_{\min}(\mathbf{H})}
  \leq \frac{\lambda \kappa(K) + \sigma_{\max}(\mathcal{D})/\sigma_{\min}(K)}{\lambda},
\]
where $\mathcal{D} = (Z \otimes K)^T SS^T (Z \otimes K)$.  Under the assumption that the
kernel is well-conditioned and the data term is dominated by $\lambda K$, $\kappa$ is moderate.
The total cost is
\begin{equation}
  O\!\bigl(n^3 + \sqrt{\kappa}\,\log(1/\varepsilon)\,(n^2 r + qr)\bigr).
  \label{eq:pcg-total}
\end{equation}

\subsection{Comparison}

\begin{center}
\begin{tabular}{lll}
  \toprule
  Method & Cost & Memory \\
  \midrule
  Direct (Algorithm~\ref{alg:direct}) & $O(q n^2 r^2 + n^3 r^3)$ & $O(n^2 r^2)$ \\
  PCG (Algorithm~\ref{alg:pcg}) & $O(n^3 + \sqrt{\kappa}\log(1/\varepsilon)(n^2 r + qr))$ & $O(n^2 + qr + nr)$ \\
  \bottomrule
\end{tabular}
\end{center}

PCG is superior when $\sqrt{\kappa}\log(1/\varepsilon) \ll nr^2$ and $q \ll n^2r$.
Under the problem assumptions $r < q$ and $n \ll M$, the PCG cost \eqref{eq:pcg-total}
is at most $O(n^3 + \sqrt{\kappa}\log(1/\varepsilon) \cdot qr)$, which can be far below
the direct method's $O(n^3 r^3)$ when $r$ is large or when many iterations are avoided
by good preconditioning.

\section{Why No $O(N)$ or $O(M)$ Operations Arise}

\begin{enumerate}[label=(\roman*)]
  \item \textbf{$(Z \otimes K)\vecop(V)$ is never formed in full.}
        The naive computation would produce a vector of length $N = nM$, costing $O(nMr)$.
        Instead, Algorithm~\ref{alg:matvec} computes only the $q$ entries corresponding
        to $\Omega$: for each $\omega$ it evaluates a single inner product of length $r$,
        after a one-time $O(n^2 r)$ preprocessing step.

  \item \textbf{The Khatri-Rao product $Z$ is accessed only via $\widetilde{Z}$.}
        $Z \in \mathbb{R}^{M \times r}$ is never formed; only the $q \times r$ submatrix
        $\widetilde{Z}$ (the rows indexed by $\{j_\omega : \omega \in \Omega\}$) is used.
        Since $q \ll M$, this avoids $O(Mr)$ storage and computation.

  \item \textbf{The MTTKRP $B = TZ$ is assumed pre-computed.}
        Standard MTTKRP algorithms for sparse tensors run in $O(qr)$ time (each nonzero
        contributes to one column of $B$), not $O(Nr)$ or $O(Mr)$.
\end{enumerate}

All operations in Algorithms~\ref{alg:matvec} and~\ref{alg:pcg} involve only:
$n \times n$ (or $n \times r$) matrices, $q$-length vectors, and $r$-dimensional inner products.
No intermediate quantity of size $N$ or $M$ is ever formed.

\end{document}
